% Auto-generated by tools/qbe.py project-article-update.
% Do not edit this file by hand; edit the proof artifacts or article sections instead.

\section{Generated Current Cycle Status}
\label{app:generated-cycle-status}

This appendix page is generated from the latest ABEIS proof cycle.  It is a
status bridge for article maintenance, not a polished theorem proof.  Stable
mathematical claims should be copied into the main text only after the
corresponding Lean declaration, source-paper anchor, and proof-note export are
available.

\paragraph{Cycle.}
Task \texttt{QBE-AUTO-002}, cycle \texttt{1}, run
\texttt{runs/20260613-182230-QBE-AUTO-002-cycle01}.  Generated at
\texttt{2026-06-13 18:33:45}.

\paragraph{Lean status signal.}
\begin{itemize}
  \item \texttt{QuantumBlockEncoding/RobinMatrix.lean:24820:  sorry}
  \item \texttt{QuantumBlockEncoding/RobinMatrix.lean:24850:  sorry}
\end{itemize}

\paragraph{Plain-language status.}
The current GHL case study is not blocked because the paper lacks a proof sketch.  It is blocked because ABEIS must turn the circuit in the paper into an exact matrix statement and prove that the clean block entry has the coefficient claimed by the paper.  In ordinary terms, the paper says that a specific path through the circuit leaves the desired coefficient, while Lean requires the corresponding matrix entry and all unwanted branches to be proved exactly.

\paragraph{Pre-Lean verifier candidates.}
These checks are necessary-condition filters, not proofs.  A failed exact
finite check usually indicates that the current Lean target, circuit
transcript, or index map is wrong.  A passing check only means the candidate
survived a cheaper diagnostic; the final claim still requires a named Lean
theorem or an explicit contract.
\begin{itemize}
  \item \textbf{active [0,0] entry.} Fast check: exact rational matrix or path-sum evaluation.  Necessary because if the exact finite entry is not the target coefficient, the Lean equality for that entry cannot be true.  Lean still must prove: a passing check is only a counterexample filter; Lean must still prove the named entry lemma.
  \item \textbf{remaining branch vanish/cancel.} Fast check: support and path checker for the remaining backend slots.  Necessary because if an unwanted clean-branch path survives numerically, the block projection cannot match the paper target.  Lean still must prove: Lean must still prove the zero/cancellation lemma in the formal circuit semantics.
  \item \textbf{Ry boundary convention.} Fast check: symbolic 2-by-2 rotation check.  Necessary because a mismatched half-angle convention changes the boundary amplitude before any Lean tactic is relevant.  Lean still must prove: Lean must still record the convention bridge as a source-supported theorem.
  \item \textbf{sparse-access map.} Fast check: finite range/injectivity/permutation check.  Necessary because a reversible oracle cannot exist for a colliding or out-of-range finite map.  Lean still must prove: Lean must still prove the reversible extension and cleanup obligations.
  \item \textbf{coefficient oracle clean branch.} Fast check: exact finite evaluator for f(x\_i)/N\_f.  Necessary because the final block entry uses this coefficient, so a wrong clean branch invalidates the target theorem.  Lean still must prove: Lean must still prove bounds, orthogonality, and unitary completion or keep them as contracts.
\end{itemize}

\paragraph{Current proof-DAG frontier.}
\begin{itemize}
  \item selected\_slot\_nonzero\_counterexample: concrete \textasciigrave{}env\textasciigrave{} with \textasciigrave{}SelectedSlot(env) = 1\textasciigrave{}; status: active obstruction leaf; Lean: proposed \textasciigrave{}oneTermRobinGamma3BoundarySelectedSlotContribution\_allOne\_nonzero\_n3\textasciigrave{}
\end{itemize}

\paragraph{Open obligation signal.}
\begin{itemize}
  \item selected-slot nonzero obstruction: Lean proposed \textasciigrave{}oneTermRobinGamma3BoundarySelectedSlotContribution\_allOne\_nonzero\_n3\textasciigrave{}; class QBE-local finite evaluator witness for the selected branch; status preferred lower2 target; open
  \item fold-forces-zero guard: Lean proposed \textasciigrave{}oneTermRobinGamma3BoundaryEvaluatedBackendFold\_forces\_selectedSlotContribution\_zero\_n3 env hFold\textasciigrave{}; class QBE-local diagnostic consequence of the compiled normal form; status fallback lower2 target; open
  \item H-free evaluated backend fold: Lean \textasciigrave{}oneTermRobinGamma3BoundaryEvaluatedBackendFoldStatement\_n3 env\textasciigrave{}; class would force \textasciigrave{}SelectedSlot(env) = 0\textasciigrave{} for all \textasciigrave{}env\textasciigrave{}; status retired as active target; \textasciigrave{}finite\_matrix\_counterexample\textasciigrave{}
  \item direct H-free selected-slot feeder: Lean proposed \textasciigrave{}oneTermRobinGamma3BoundaryActiveSelectedSlotEvalFeeder\_n3 env\textasciigrave{}; class active row \textasciigrave{}0\textasciigrave{} versus selected slot \textasciigrave{}2\textasciigrave{} / full index \textasciigrave{}32\textasciigrave{}; status retired; \textasciigrave{}shape\_or\_register\_gap\textasciigrave{}
  \item source-prepared recovery: Lean \textasciigrave{}oneTermRobinGamma3BoundarySourcePreparedProjectionTarget\_activeEval\_of\_evaluatedBackendFold\_n3 H env hUniform hFold\textasciigrave{}; class compiled downstream route under explicit clean-column contract; status still valid only after a corrected valid fold theorem
  \item source-prepared retarget: Lean corrected source/register target after obstruction; class source-contract audit plus QBE-local finite matrix semantics; status blocked until lower2 obstruction or reviewer restatement
\end{itemize}

\paragraph{Open GHL contribution obligations.}
\begin{itemize}
  \item \texttt{Uindic}: Indicator unitary; status Permutation and self-inverse helpers compile, but this is not yet the final block-encoding proof..
  \item \texttt{RyBoundary}: Boundary-controlled Ry rotations; status Active convention audit: the Ry angle convention must be source-supported before closing the boundary amplitude proof..
  \item \texttt{RobinTheorem}: One-term Robin block-encoding theorem; status Main active target: the theorem-facing Lean statement is not closed..
  \item \texttt{GammaSlices}: Gamma wavefunction slices and coefficient product; status Several finite boundary lemmas exist; the final projection/product bridge is still open..
  \item \texttt{FigRobin}: Figure 4 Robin circuit transcript; status Transcript guard compiles for visible indicator dagger and H preparation; the active seven-gate backend remains a component..
  \item \texttt{OneD}: One-dimensional Hamiltonian block encoding; status Deferred until the one-term Robin theorem closes..
  \item \texttt{MultiD}: Multidimensional block encoding; status Planned; depends on one-term and LCU abstractions..
\end{itemize}

\paragraph{Open external technical lemma obligations.}
\begin{itemize}
  \item \texttt{tl-ghl-lemma1-banded-sparse-access}: contract-only; next action Keep as explicit external contract until the exact cited construction is formalized or imported as a theorem..
  \item \texttt{tl-ghl-lemma3-sparse-amplitude}: contract-only; next action Maintain the clean-branch contract and isolate any missing unitarity proof as an external technical lemma..
  \item \texttt{tl-ghl-theorem5-piecewise-polynomial-of}: contract-only; next action Use only as a named contract in the current theorem; do not invent stronger smoothness or range assumptions..
  \item \texttt{tl-uniform-sparse-register-preparation}: contract-only; next action Keep as a theorem-facing contract unless the exact state-preparation proof is needed to close a resource theorem..
  \item \texttt{tl-ry-boundary-amplitude-convention}: obligation; next action Do not silently change the paper angle; prove the convention bridge or record the exact cited convention..
  \item \texttt{tl-qsvt-blockencoding-simulation}: paper-cited; next action Defer until the one-term Robin theorem and LCU combination are closed..
  \item \texttt{tl-clean-block-definition}: contract-only; next action Use as the final target predicate; ensure circuit entry lemmas are bridged to this semantic definition..
\end{itemize}

\paragraph{Recent typed verifier feedback.}
\begin{itemize}
  \item Leaf \texttt{branch\_correct\_evaluated\_backend\_fold\_route\_guard}, class \texttt{finite\_matrix\_counterexample}; next route: Prove the necessary-condition route guard or formalize the selected-slot nonzero counterexample, then middle should retire or restate the H-free evaluated fold. Source-prepared recovery remains valid only under explicit hUniform after a valid fold theorem..
  \item Leaf \texttt{branch\_correct\_evaluated\_backend\_fold}, class \texttt{finite\_matrix\_counterexample}; next route: Do not assign lower2 to prove oneTermRobinGamma3BoundaryEvaluatedBackendFoldStatement\_n3 env as an all-env theorem in the current row-0/slot-2 finite witness. Middle should repair the source contract or proof-DAG leaf: either change the active row/register target to the branch-correct selected source entry, add an explicit source-backed coefficient/vanish obligation for selectedSlotContribution = 0 if that is truly intended, or record a source\_translation\_gap before further Lean proof search..
  \item Leaf \texttt{branch\_correct\_evaluated\_backend\_fold\_selected\_slot\_zero\_normal\_form}, class \texttt{symbolic\_bridge\_gap}; next route: prove selected-slot scalar eval zero or classify it as the next source/register obstruction.
  \item Leaf \texttt{branch\_correct\_evaluated\_backend\_fold}, class \texttt{symbolic\_bridge\_gap}; next route: Lower2 should prove exactly one evalWith-level branch-correct full backend-fold leaf, preferably ActiveEntry(env) = BackendFold(env), without adding hUniform or reviving the selected-slot shortcut..
  \item Leaf \texttt{source\_prepared\_finite\_composition\_leaf}, class \texttt{symbolic\_bridge\_gap}; next route: prove the unwrapped sparse-clean evaluated equality or a source-backed finite-composition lemma feeding it; do not revive the H-free row-0 to slot-2 feeder.
  \item Leaf \texttt{source\_prepared\_finite\_composition\_leaf}, class \texttt{symbolic\_bridge\_gap}; next route: Lower2 should prove exactly one source-prepared active/prepared finite-composition leaf, preferably oneTermRobinGamma3BoundaryUncastActivePreparedCompositeEvalStatement\_n3 H env or (oneTermRobinGamma3BoundaryActivePreparedEntryTarget\_n3 H).entryEqualityStatement. Do not add Uniform(H) to the active/prepared field and do not revive the H-free row-0 to selected slot-2 feeder..
\end{itemize}

\paragraph{Recent proof-attempt memory.}
\begin{itemize}
  \item \texttt{proof-attempts/QBE-AUTO-002/source-prepared-backend-fold-obstruction-middle-packet-20260613-182230.md}
  \item \texttt{proof-attempts/QBE-AUTO-002/source-prepared-backend-fold-lower1-route-guard-20260613-181816.md}
  \item \texttt{proof-attempts/QBE-AUTO-002/source-prepared-backend-fold-lower3-necessary-condition-20260613-181816.md}
  \item \texttt{proof-attempts/QBE-AUTO-002/source-prepared-backend-fold-lower2-20260613-181849.md}
  \item \texttt{proof-attempts/QBE-AUTO-002/source-prepared-backend-fold-middle-packet-20260613-180059.md}
\end{itemize}

\paragraph{Article-writing rule.}
The project report may discuss this cycle as process evidence.  It must not
claim completion of the first case study \citep{guseynovHuangLiu2025} unless
the final theorem-facing block-extraction statement is closed in Lean and the
Markdown/LaTeX proof map has been synchronized.
